\scanpage{201}
$\mathbf C$, i.e. $\overline{k(a)}\simeq\mathbf R$, \inserted{isom uni. can.} alors $\xi_{x,\mathbf Z}$ est encore défini par la valuation archimédienne sur $k(a)$, mais \deleted{\unsure{suivant que}} \inserted{\unsure{compte tenu|indépendant}} des choix possibles dans la clôture algébrique de $k(a)$, le groupe de Galois de $\mathbf C=\overline{\mathbf R}$ sur $\mathbf R$ opère dans $\xi_{x,\mathbf Z}$, \unsure{notons} $\tau$ \uline{cette opération}, qui est involutive. \inserted{on écrit aussi $\xi_{a,\mathbf Z}$ au lieu de $\xi_{x,\mathbf Z}$,} [On ne la confondra pas avec l'endom \uline{$\frac12$ linéaire} de $\xi_{\bar a,\infty}$ correspondant à $\tau$ ; il n'y est pas question par \unsure{exple|exemple} d'un réseau stable].

\margin{\unsure{a devient le} \uline{\unsure{plus}} \unsure{défini que pt à val. dans} $\mathbf C$}

Dans ce cas, les invariants sont
\diagram{\begin{tikzcd}[column sep=large,row sep=small]
\xi_{a,\infty}\in T_\alpha(\mathbf R)\arrow[dr,no head] & \\
& \xi_{\bar a,\infty}\in T_\alpha(\mathbf C)\\
\begin{gathered}\deleted{\xi_{a,\infty}\otimes_{\mathbf R}\mathbf C\simeq}\\\xi_{\bar a,\mathbf Z}\in T_\alpha(\mathbf Z)\end{gathered}\arrow[ur,no head]&
\end{tikzcd}}
\[
\xi_{\bar a,\infty}\overset{\text{\uuline{can}}}{\simeq}\xi_{a,\infty}\otimes_{\mathbf R}\mathbf C\simeq\xi_{\bar a,\mathbf Z}\otimes_{\mathbf Z}\mathbf C
\]

$\pi_1(X)$ n'opère pas à priori dans $\xi_{a,\mathbf Z}$, les seules opérations sont celles de $\tau$ si $a$ est \uline{réel}. Cependant, sur $\xi_{\bar a,\infty}$, il y a (en plus éventuellement de $\tau$, si $a$ est réel) les opérations de $\operatorname{Aut}_{k(a)}(\unsure{k})$, qui sont $\frac12$ linéaires.

\mspage{200}
